In this section, we detail the \textit{course phase} of the proposal, during
which we developed two semantic extensions that enrich the reasoning
capabilities of classical e-graphs: \textit{disequality reasoning} and 
\textit{conditional reasoning}.

\subsection{Disequality Reasoning}
First, I contributed to \textit{dis/equality graphs}, which extend
e-graphs with direct support for disequality~\cite{zakhour2025disequality}.
Equality is represented as usual, while a disequality is stored as a
first-class, undirected \textit{disequality edge} between two e-classes. This
contrasts with the \textit{embedding} techniques prevalent in SMT solvers and
verifiers, which reify (dis)equality into the term language, using a symbol such
as \texttt{eq} or \texttt{ne} and saturating it against \texttt{false}.
Embeddings enlarge the e-graph with extra e-nodes and e-classes, or require
saturation; a disequality edge only requires storing a pair of e-classes plus
some ad-hoc bookkeeping without affecting the term space, reducing both time and
memory overhead.

We studied this design formally and empirically. The work introduced a formal
framework for reasoning about e-graphs and used it to prove that an e-graph
coincides with the reflexive, symmetric, transitive, and
congruent closure of the asserted equalities. Generalizing to dis/equality
graphs, we proved that they detect any contradiction between the recorded
equalities and the asserted disequalities, and we gave an analytical bound
showing them to be asymptotically more efficient than the embeddings.
We implemented them as an
extension to egg~\cite{willsey2021egg} and evaluated them in an SMT solver and
the Propel automated theorem prover~\cite{zakhour2024algebraic, zakhour2023crdt} on
standard benchmarks, where direct support for disequalities outperformed 
embedding-based encodings, in line with the analytical result.

\subsection{Conditional Reasoning}
Second, I authored \textit{versioned e-graphs}, which extend e-graphs to
represent equalities that hold only conditionally~\cite{cesario2026versioned}.
Under proof by cases, and branching settings more generally, each branch assumes
a different set of equalities on top of a shared base. A traditional e-graph
encodes a single global equality relation, so these contextual equalities are
handled by keeping one e-graph per branch, which replicates the equalities the
branches share. A versioned e-graph instead encodes many equivalence sets, one
per \textit{version}, in a single structure that maximizes sharing across them,
resting on a \textit{versioned union-find} that records the equalities common to
several versions once.

As before, we studied this design formally and empirically. The work formalized
versioned union-finds and versioned e-graphs and proved them correct, and
analyzed the amortized time and space complexity of the versioned union-find
against the cloning of a separate e-graph per branch. We realized the design in
Veg, a standalone Rust library inspired by egg~\cite{willsey2021egg} and 
evaluated it against cloning, a variant based on persistent data structures, and
the state of the art~\cite{singher2024easteregg} across two case studies, an EUF
solver and the TheSy theorem prover~\cite{singher2021thesy}, where versioned
e-graphs used up to $30\%$ less memory and ran up to $4\times$ faster, with the
gap widening as the number of explored terms and branches grew.